\section{Introduction}\label{sec:intro}
Interactive educational systems that combine automated scoring, language models, and reinforcement-learning (RL) policies are increasingly described as multi-agent systems~\cite{autogen,metagpt}, even when a single monolithic controller executes every component in a fixed procedural pipeline. Whether the multi-agent label reflects a genuine scientific distinction depends on how authority and state are distributed~\cite{wj95}. If components operate as specialized autonomous agents, critical questions arise regarding coordination, decision consistency, and architectural necessity: how do specialized agents coordinate their recommendations, how does their interaction affect session trajectories, and does partitioning responsibilities across specialized agents provide measurable value over a unified controller?

We investigate these questions within an orchestrated spoken technical-interview preparation system whose deployed architecture we audit rather than assume. In its baseline configuration, the system executes an automated interview workflow driven by a single learned policy (a frozen Proximal Policy Optimization controller~\cite{ppo,sb3}) operating over a 6-dimensional state vector without concept memory. Every other pedagogical choice is governed by static threshold rules in the central hub. We extend this system by introducing two specialized, stateful agents behind a feature flag, leaving the underlying evaluator, frozen PPO policy, and candidate contract intact:
\begin{enumerate}
\item \textbf{Strategy Agent:} Maintains a private weakness ledger across turns, tracking persistent conceptual deficiencies and score trajectories to propose dynamic difficulty transitions and proactive follow-up probes.
\item \textbf{Coaching Agent:} Maintains a private coaching history ledger, tracking previously delivered hints and examples to propose feedback depth, instructional emphasis, and targeted remediation tags.
\end{enumerate}
The agents possess strictly bounded authority: they only \emph{propose}. A deterministic coordinator validates all proposals against formal schemas and authority matrices, an arbitrator resolves competing recommendations against safety bounds and the PPO policy, and the central hub remains the sole system of record. Crucially, the Strategy Agent communicates pedagogical intent to the Coaching Agent via unidirectional peer messaging, enabling aligned instructional guidance without coupling private state ledgers.

This architecture enables a rigorous empirical evaluation of multi-agent coordination and specialization. Rather than evaluating generic agent reliability or treating multi-agent design as an axiomatic virtue, we formulate three precise research questions:
\begin{description}
\item[RQ1 (Coordination and intervention dynamics):] How does stateful coordination between the Strategy and Coaching agents affect the decisions and session trajectories of the interview preparation system, compared with the frozen single-policy baseline ($B_0$)?
\item[RQ2 (Empirical effects of agent specialization):] Does separating Strategy and Coaching into distinct specialized agents yield measurable decision divergence or operational trade-offs compared to a unified technical coaching agent operating on equivalent information and tasks?
\item[RQ3 (Coordination mechanisms):] Which implemented coordination components---stateful memory, peer messaging, confidence threshold gating, or PPO interaction---are responsible for observable differences in system decisions?
\end{description}

\noindent\textbf{Summary of Findings and Contributions.}
\begin{itemize}
\item \emph{Coordination Dynamics (RQ1):} In deterministic replay across 160 sessions (operating on a fixed set of 64 technical-interview replay instances across 20 question permutation orders), stateful coordination alters difficulty trajectories in all 160 sessions relative to the frozen baseline. The Strategy Agent smooths trajectory oscillations, reducing mean level changes from 15.12 to 12.45 per session ($-2.67$, cluster-bootstrap 95\% CI \phchgci{}) and mean reversals from 11.74 to 9.87 ($-1.88$, cluster CI \phrevci{}). Agreement with the uncoordinated PPO action decreases from 1.000 to 0.774 ($-0.226$, cluster CI \phagrci{}). Proactive probes introduce an average of $+0.59$ additional turns per session (cluster CI \phfuci{}). However, a cluster sensitivity analysis reveals that this probing increase is driven entirely by a single replay cluster ($k=1$); omitting cluster 1 yields $\Delta = +0.00$ extra turns.
\item \emph{Limits of Specialization (RQ2):} Under the tested deterministic replay configuration, the split and unified architectures produced identical decisions and difficulty trajectories across all 160 sessions (1,280 turns; $\mathbb{P}(\text{agree}) = 1.0$), demonstrating that architectural decomposition alone did not alter the implemented decision policy. Ablating the Coaching Agent ($M_3$) reproduces the full two-agent system's difficulty trajectory in 160 of 160 sessions, as difficulty adaptation is structurally independent of coaching state. Specialization introduces measurable end-to-end orchestration overhead ($+1.04$\,ms median turn latency, 47.7 agent calls per session) alongside micro-benchmark serialization costs ($+0.0054$\,ms per turn; $0.0092$\,ms vs $0.0038$\,ms). Under the tested configuration, specialization primarily manifests as an architectural decomposition rather than a change in the resulting decision policy, establishing distinct state ledgers and isolated decision records rather than measured algorithmic gains.
\item \emph{Mechanism Decomposition (RQ3):} Through controlled ablations and a multi-level threshold sweep ($\tau \in [0.0, 0.9]$), we isolate the confidence threshold as the critical gating mechanism. Setting $\tau = 0.0$ collapses level changes to 2.225 and PPO agreement to 0.014, identical to removing arbitration entirely ($M_4$). Peer messaging ($M_5$) alters coaching depth and tone without destabilizing difficulty. Disabling PPO ($M_7$) alters trajectories in only 4 of 160 sessions (2.5\%), contradicting the registered prediction of $\ge 50\%$ divergence and revealing strong alignment between the Strategy Agent heuristic and the baseline rule fallback.
\end{itemize}

\noindent\textbf{What Kind of Contribution This Is.} The specialized agents, typed contracts, and authority matrix constitute \emph{architecture}; the peer messaging and threshold gating constitute \emph{coordination mechanisms}; the replay ablations and cluster-bootstrap evaluations provide \emph{empirical evidence on coordination and the limits of specialization}.

\noindent\textbf{Explicit Disclaimers.} All replay experiments execute deterministically with recorded evaluator outputs (`EVALUATOR\_MOCK\_MODE=1`); no live large language model was invoked during replay. We make \textbf{no claims} regarding candidate learning, pedagogical efficacy, interview placement success, or LLM hallucination rates. Fewer trajectory reversals reflect behavioral stability, not proven pedagogical superiority.
